\section{About the Committee}

\subsection{Mandate}

The World Bank is the world’s largest development institution, providing over US\$100 billion per year through loans, grants, and guarantees.\footnote{World Bank. (2025a). Annual report 2025. World Bank Group. \url{https://www.worldbank.org/en/about/annual-report}} It was established alongside the International Monetary Fund (IMF) at the 1944 Bretton Woods Conference, and before turning its attention to the developing world, its initial purpose was to finance the reconstruction of Europe after the Second World War.\footnote{World Bank. (n.d.-a). International Bank for Reconstruction and Development. \url{https://www.worldbank.org/en/who-we-are/ibrd}} Today, the Bank has two targets: to end extreme poverty and to boost shared prosperity. In contrast to the IMF (which was created to safeguard the stability of the global monetary and financial system), the World Bank was created to finance long-term development. Such as roads, energy supply, schools, hospitals, and institutions that facilitate growth in poorer countries.

Apart from its loan book, the Bank has a wide range of tools that together define its influence. Its financing typically comes with conditions: a borrowing government may be asked to meet requirements on transparency, fiscal management, or the protection of minimum levels of social spending as part of the arrangement.\footnote{International Development Association. (2023). Debt. World Bank. \url{https://ida.worldbank.org/en/financing/debt}} Beyond money, the Bank is one of the world’s most important producers of economic data and analytical frameworks that governments, markets, and other institutions rely on to assess risk and guide policy. Additionally, the Bank acts as a convener, bringing together parties, including governments, creditors, and international organisations, to address challenges that no single actor can resolve alone.

\subsection{Structure: The World Bank Group}

It is important to be precise about what “the World Bank” means. Strictly, the institution is the World Bank Group (WBG), a family of five organisations.\footnote{World Bank Group. (n.d.). Who we are. \url{https://www.worldbank.org/en/who-we-are}} Two of them are responsible for lending to governments and are together often referred to as “the World Bank”:

The International Bank for Reconstruction and Development (IBRD) lends to middle-income and creditworthy lower-income governments. It raises the vast majority of its money by issuing bonds on international capital markets. Because it holds a triple-A credit rating, it can borrow cheaply and pass these favourable terms on to borrowers.\footnote{World Bank. (n.d.-a). International Bank for Reconstruction and Development. \url{https://www.worldbank.org/en/who-we-are/ibrd}} IBRD loans are offered at close to market rates, similar to the terms a well-rated sovereign borrower would receive.

The International Development Association (IDA) is the Bank’s concessional arm, meaning it lends on far softer terms than any commercial lender would offer, and increasingly provides grants to the world’s poorest countries.\footnote{International Development Association. (2024). History: What is IDA? World Bank. \url{https://ida.worldbank.org/en/about/history}} A country generally becomes IDA-eligible when its gross national income (GNI) per capita falls below US\$1,325, and/or the country lacks creditworthiness to borrow from the IBRD.\footnote{International Development Association. (2026). Financing. World Bank. \url{https://ida.worldbank.org/en/financing}} Unlike the IBRD, the IDA is funded largely by contributions (“replenishments”) negotiated every three years by donor governments.\footnote{Congressional Research Service. (2025). The World Bank (IF11361). \url{https://www.congress.gov/crs-product/IF11361}}

Three other institutions complete the Group and are worth knowing to understand the World Bank Group and its scope as a whole. The International Finance Corporation (IFC) is the Group’s private-sector arm that invests directly in private companies in developing countries to stimulate growth and create jobs. The Multilateral Investment Guarantee Agency (MIGA) encourages investment into these economies by insuring private investors against political risks such as expropriation, conflict, or a government blocking the transfer of currency abroad. The International Centre for Settlement of Investment Disputes (ICSID) serves as a forum for solving legal disputes between foreign investors and host governments. For the purposes of this committee, the two actors that matter most are the IBRD and, above all, IDA, since their lending to national governments is what this committee’s topic is all about. Although the IFC, MIGA, and ICSID work in a different lane, delegates looking for creative approaches may still find that the Group’s wider toolkit has something to offer for their solutions.

\subsection{Governance and Decision-Making}

At IEUMUN, the World Bank will run like any other General Assembly committee. Each delegation holds one equal vote. Debate runs through moderated and unmoderated caucuses, and delegates work together toward a resolution, which is debated upon and adopted by majority vote. At the real World Bank, votes are weighted by financial contribution, which is a major reason reform can be particularly difficult. The most fundamental decisions require an 85\% supermajority, and the US, for example, holds roughly 16\% of the vote, effectively giving it veto power.\footnote{Congressional Research Service. (2025). The World Bank (IF11361). \url{https://www.congress.gov/crs-product/IF11361}} Our committee sets that feature aside, but delegates should keep in mind that in reality, the Bank’s largest shareholders cannot simply be outvoted.
